% ===== PRÉAMBULE CANONIQUE — à coller VERBATIM en tête de CHAQUE .tex =====
\documentclass[11pt,a4paper]{article}
\usepackage[utf8]{inputenc}\usepackage[T1]{fontenc}\usepackage{lmodern}
\usepackage[french]{babel}
\usepackage{amsmath,amssymb,mathtools}\usepackage{icomma}
\usepackage{siunitx}\sisetup{locale=FR}  % \qty{}{}, \si{}, \ang{}
\usepackage{xcolor}\usepackage{colortbl}
\usepackage{tikz}\usepackage{pgfplots}\pgfplotsset{compat=1.18}
\usepackage[siunitx,europeanresistors]{circuitikz} % circuits (élec.)
\usepackage[most]{tcolorbox}\usepackage{enumitem}\usepackage{array,booktabs}
\usepackage[a4paper,margin=2.2cm]{geometry}
\usetikzlibrary{arrows.meta,calc,angles,quotes,positioning,patterns,%
  backgrounds,shapes.geometric,decorations.pathmorphing,decorations.markings,intersections}
\definecolor{primary}{RGB}{222,49,99}\definecolor{primarydark}{RGB}{150,22,55}
\definecolor{defcol}{RGB}{0,114,178}\definecolor{methcol}{RGB}{52,92,148}
\definecolor{excol}{RGB}{0,135,90}\definecolor{attncol}{RGB}{200,40,55}
\tcbset{boxbase/.style={enhanced,breakable,boxrule=0.9pt,arc=2.5pt,
  left=3mm,right=3mm,top=2.3mm,bottom=2.3mm,fonttitle=\bfseries,coltitle=white}}
% --- boîtes TUTORIEL (noms EXACTS, ne pas renommer) ---
\newtcolorbox{cle}[1][]{boxbase,colback=defcol!6,colframe=defcol,
  title={\textsc{À retenir}\ifx\relax#1\relax\relax\else\textmd{ : \itshape #1}\fi}}
\newtcolorbox{methode}[1][]{boxbase,colback=methcol!7,colframe=methcol,sharp corners,
  title={\textsc{Méthode}\ifx\relax#1\relax\relax\else\textmd{ --- \itshape #1}\fi}}
\newtcolorbox{exemplebox}[1][]{boxbase,colback=excol!5,colframe=excol,
  title={\textsc{Exemple}\ifx\relax#1\relax\relax\else\textmd{ : \itshape #1}\fi}}
\newtcolorbox{piege}[1][]{boxbase,colback=attncol!6,colframe=attncol,
  title={\textsc{Piège}\ifx\relax#1\relax\relax\else\textmd{ : \itshape #1}\fi}}
\newtcolorbox{remarque}[1][]{boxbase,colback=black!4,colframe=black!45,
  title={\textsc{Remarque}\ifx\relax#1\relax\relax\else\textmd{ : \itshape #1}\fi}}
% --- boîtes EXERCICE / CORRIGÉ (noms EXACTS, auto-comptées) ---
\newtcolorbox[auto counter]{exo}[1][]{boxbase,colback=primary!4,colframe=primary,
  title={\textsc{Exercice}~\thetcbcounter\ifx\relax#1\relax\relax\else\textmd{ --- \itshape #1}\fi}}
\newtcolorbox[auto counter]{corrige}[1][]{boxbase,colback=excol!4,colframe=excol,
  title={\textsc{Corrigé}~\thetcbcounter\ifx\relax#1\relax\relax\else\textmd{ --- \itshape #1}\fi}}
\newcommand{\vv}[1]{\vec{#1}}\newcommand{\ds}{\displaystyle}
\newcommand{\R}{\mathbb{R}}\newcommand{\N}{\mathbb{N}}\newcommand{\Z}{\mathbb{Z}}
% =========================================================================
\begin{document}
\usetikzlibrary{babel}
\begin{corrige}[bloc parallèle en série : retrouver $R$]
\begin{enumerate}[label=\alph*)]
  \item Bloc parallèle (3 résistances égales) : $R_{\parallel}=\dfrac{12}{3}=\qty{4}{\ohm}$.
        La puissance totale donne la résistance équivalente vue du générateur :
        \[ P=\frac{U^{2}}{R_{\text{eq}}}\;\Rightarrow\; R_{\text{eq}}=\frac{U^{2}}{P}=\frac{20^{2}}{40}=\qty{10}{\ohm}. \]
  \item $R_{\text{eq}}=R+R_{\parallel}$, donc $R=R_{\text{eq}}-R_{\parallel}=10-4=\qty{6}{\ohm}$.
\end{enumerate}
\end{corrige}

\begin{corrige}[réduire un réseau mixte pour trouver $I$]
Bloc parallèle : $\qty{30}{\ohm}\,\|\,\qty{30}{\ohm}=\dfrac{30}{2}=\qty{15}{\ohm}$. Il est en série avec
$\qty{5}{\ohm}$ et $\qty{4}{\ohm}$ :
\[ R_{\text{eq}}=5+15+4=\qty{24}{\ohm},\qquad I=\frac{U}{R_{\text{eq}}}=\frac{24}{24}=\qty{1}{\ampere}. \]
\end{corrige}

\begin{corrige}[diviseur de courant et puissance d'une branche]
\begin{enumerate}[label=\alph*)]
  \item Les deux résistances ont la même tension : $U=\qty{2}{\ohm}\times I_{2\Omega}=2\times1=\qty{2}{\volt}$.
  \item Dans la résistance de $\qty{4}{\ohm}$ : $P=\dfrac{U^{2}}{R}=\dfrac{2^{2}}{4}=\qty{1}{\watt}$
        (l'intensité y vaut $I_{4\Omega}=U/4=\qty{0.5}{\ampere}$).
\end{enumerate}
\end{corrige}

\begin{corrige}[remplacer une résistance]
\begin{enumerate}[label=\alph*)]
  \item $R_{\text{eq}}=\dfrac{30}{3}=\qty{10}{\ohm}$ ; la tension du générateur est
        $U=R_{\text{eq}}\,I=10\times2=\qty{20}{\volt}$.
  \item Nouvelle résistance équivalente :
        \[ \frac{1}{R_{\text{eq}}'}=\frac{1}{30}+\frac{1}{30}+\frac{1}{7{,}5}=\frac{1}{15}+\frac{1}{7{,}5}
           =\frac{1}{5}\;\Rightarrow\; R_{\text{eq}}'=\qty{5}{\ohm}. \]
        À tension fixe, $I'=\dfrac{U}{R_{\text{eq}}'}=\dfrac{20}{5}=\qty{4}{\ampere}$ : le courant a
        \textbf{doublé} (la $R_{\text{eq}}$ a diminué).
\end{enumerate}
\end{corrige}

\begin{corrige}[pile réelle : à vide, en charge, en court-circuit]
On utilise $U=E-rI$.
\begin{enumerate}[label=\alph*)]
  \item À vide, $I=0$ : $U=E=\qty{9}{\volt}$.
  \item En charge : $U=9-1\times2=\qty{7}{\volt}$.
  \item Court-circuit ($U\approx0$) : $I_{\text{cc}}=\dfrac{E}{r}=\dfrac{9}{1}=\qty{9}{\ampere}$.
  \item $P=r\,I_{\text{cc}}^{2}=1\times9^{2}=\qty{81}{\watt}$ : c'est pourquoi une pile en court-circuit
        chauffe dangereusement.
\end{enumerate}
\end{corrige}

\end{document}
